\section{Related Work}
\label{sec:related_work}
\textbf{LLMs for robot design.} Using language models to assist robot design has emerged only recently~\cite{qiu2026robomorph, ringel2025text2robot}. RoboMorph~\cite{qiu2026robomorph} couples an LLM with evolutionary search to generate modular morphologies from language specifications. These pipelines typically optimize morphology under a fixed or hand-tuned objective and do not jointly reason about reward shaping for changed bodies. \abbrev{} targets design--reward co-design: morphology edits are paired with reward code and validated through simulator-in-the-loop training.

\textbf{LLMs for reward design.} A parallel thread of research explores using LLMs to automate reward design~\cite{ma2024eureka, kwon2023reward, yu2023language, xie2024text2reward}. For example, Eureka~\cite{ma2024eureka} demonstrated that LLM-generated reward code can outperform hand-engineered rewards on challenging tasks by leveraging domain knowledge. Most existing approaches optimize reward code for a fixed morphology. \abbrev{} conditions reward generation on the proposed morphology and evaluates candidates in simulation, enabling reward shaping that co-adapts with morphological changes.

\textbf{Co-optimization of morphology and policy.} Given the strong coupling between a robot's body and its controller~\cite{pfeifer2006morphological, sims1994evolving, cheney2018scalable}, prior work has studied joint optimization of morphology and control. Reward hacking and specification gaming motivate our separation of the learned reward $r$ (used for training) from the standardized evaluation score $S$ (used for selection)~\cite{amodei2016concrete, krakovna2020specification}. RoboMoRe~\cite{fang2025robomore} alternates between LLM-proposed morphology edits and reward hints in a coarse-to-fine fashion with a single LLM alternating roles. \abbrev{} differs in three ways: (1) it explicitly separates design and reward generation across two specialized agents, (2) it uses a dialectical thesis--antithesis--synthesis structure where agents critique each other's proposals before refinement, and (3) it employs criterion-specific judges that provide multi-objective feedback rather than relying on a single unified objective.

\textbf{Quality-diversity methods.} Quality-diversity (QD) algorithms such as MAP-Elites~\cite{mouret2015illuminating} maintain archives of diverse, high-performing solutions across a behavior space~\cite{pugh2016quality}. These methods excel at discovering broad repertoires of designs but require hand-specified behavior descriptors and do not incorporate semantic reasoning about \emph{why} certain designs succeed or fail. \abbrev{} has a different goal. Instead of filling a behavior archive, it seeks a single high-performing design through LLM-driven critique and refinement. The approaches are complementary---combining debate-generated candidates with QD archives is a promising direction we leave to future work.

\textbf{Multi-agent LLM debates.} Multi-agent debate has been studied as a mechanism to improve reasoning and evaluation~\cite{irving2018ai, liang2023encouraging, du2024improving, chan2023chateval, gu2024survey}. For example, \cite{chan2023chateval} show that debate-based protocols can increase evaluation reliability, and mixtures of judges have been used to reduce reward hacking in preference learning~\cite{xu2024moj}. \abbrev{} adapts this idea to robotics by grounding debate in a physics simulator. LLM judges condition on measured metrics and provide textual critiques, while selection is based solely on the task score $S$. Unlike prior debate frameworks that rely on LLM-learned judges or preference aggregation, \abbrev{} uses physics metrics as ground truth and LLM judges for qualitative guidance only. To our knowledge, \abbrev{} is the first to combine simulator-grounded multi-agent debate with criterion-specific LLM critics for morphology--reward co-design.

\begin{table}[H]
\centering
\caption{Comparison with related LLM-based robotics methods. \cmark{} = directly optimized or used as a primary mechanism, \xmark{} = not a primary target. \textbf{Takeaway:} \abbrev{} combines morphology search, reward-code search, iterative refinement, and role-separated multi-agent debate.}
\label{tab:related_work_comparison}
\small
\setlength{\tabcolsep}{3.5pt}
\begin{tabular}{@{}lcccc@{}}
\toprule
\textbf{Method} & \textbf{Morph.} & \textbf{Reward} & \textbf{Iter.} & \textbf{Role sep.} \\
\midrule
RoboMorph & \cmark & \xmark & \cmark & \xmark \\
Eureka & \xmark & \cmark & \cmark & \xmark \\
Text2Robot & \cmark & \xmark & \xmark & \xmark \\
VLMgineer & \cmark & \cmark & \cmark & \xmark \\
RoboMoRe & \cmark & \cmark & \cmark & \xmark \\
\rowcolor{gray!15}
\textbf{\abbrev{} (ours)} & \cmark & \cmark & \cmark & \cmark \\
\bottomrule
\end{tabular}
\end{table}

\textbf{Summary of key differences.} Table~\ref{tab:related_work_comparison} summarizes the comparison axes most relevant to \abbrev{}. The closest methods are VLMgineer~\cite{gao2025vlmgineer} and RoboMoRe~\cite{fang2025robomore}, which also pair design changes with control objectives, but they do not use role-separated design--control debate with criterion-specific judges. \abbrev{} therefore differs in protocol rather than only in objective: specialized agents propose and critique design--reward pairs, while physics metrics ground both feedback and final archive selection.
